\chapter{Полный калибровочно-взвешенный носитель рёбер}
\label{ch:v7-full-gauge-weighted-edge-carrier-gate}

\section{Минимальный подъём полных блоков}

Редуцированный Hodge-родитель использовал одно комплексное число на ребро.
Теперь раскроем минимальный калибровочный мультиплет каждого блока, стягивая
тождественным интертвинером общий фактор концов. Для представления
\begin{equation}
 R_e=(d_3,d_2)_{Y_e}
 \label{eq:v7-full-edge-representation}
\end{equation}
следовые индексы в стандартной нормировке равны
\begin{equation}
 I_1(e)=d_3d_2Y_e^2,qquad
 I_2(e)=d_3T_2(R_e),qquad
 I_3(e)=d_2T_3(R_e),
 \label{eq:v7-full-edge-indices}
\end{equation}
где $T(\mathbf2)=T(\mathbf3)=1/2$. Именно представления на конечном
гильбертовом пространстве определяют коэффициенты gauge-кривизны в
спектральном действии \cite{v7-full-edge-spectral-action}.

Полный список одиннадцати блоков имеет вид
\begin{equation}
\begin{array}{c|c|c}
e&R_e&\text{статус}\\ \hline
L_LY_R&(1,1)_0&*\\
Q_LY_R&(3,1)_{2/3}&*\\
X_LX_R&(1,1)_0&*\\
X_Le_R&(1,1)_0&*\\
X_Lu_R&(3,1)_{5/3}&*\\
Y_LY_R&(1,1)_0&*\\ \hline
L_LX_R&(1,2)_{1/2}&+\\
X_LY_R&(1,2)_{1/2}&+\\
X_Ld_R&(3,1)_{2/3}&+\\
Y_LX_R&(1,2)_{1/2}&+\\
Y_Le_R&(1,2)_{1/2}&+
\end{array}
 \label{eq:v7-full-edge-multiplet-table}
\end{equation}
Звёздочкой отмечены шесть рёбер, выбранных $P_*$; знаком плюс --- пять
конкурирующих. Стрелки диаграммы кодируют массовые и Yukawa-блоки между
хиральными мультиплетами, поэтому такое раскрытие является необходимым до
физической интерпретации потенциала \cite{v7-full-edge-krajewski}.

\section{Что переживает физические веса}

Пусть каждый полный блок имеет положительный следовой вес $w_e$. Тогда
естественный блочно-радиальный подъём прежнего потенциала равен
\begin{equation}
 \mathcal S_{\rm full}(\phi)
 =\sum_{e\in E_*}w_e(\|\phi_e\|^2-\mu^2)^2
 +\sum_{e\notin E_*}w_e
 (\|\phi_e\|^4+2\mu^2\|\phi_e\|^2).
 \label{eq:v7-full-edge-weighted-potential}
\end{equation}
При любых $w_e>0$ множество минимумов сохраняет опору
\begin{equation}
 \|\phi_e\|=\mu\quad(e\in E_*),qquad
 \phi_e=0\quad(e\notin E_*).
 \label{eq:v7-full-edge-weighted-vacuum}
\end{equation}
Следовательно, положительные физические кратности не меняют сам селектор
``шесть из одиннадцати''.

Суммарные комплексные размерности выбранной и конкурирующей частей равны
\begin{equation}
 d_* =10,qquad d_+=11.
 \label{eq:v7-full-edge-dimensions}
\end{equation}
До gauge-quotient гессиан блочно-радиального потенциала имеет сигнатуры
\begin{equation}
 (n_-,n_0,n_+)_{0}=(20,0,22),qquad
 (n_-,n_0,n_+)_{*}=(0,14,28).
 \label{eq:v7-full-edge-hessian-signatures}
\end{equation}
Четырнадцать нулей являются угловыми направлениями сфер ненулевых полных
мультиплетов; часть из них может быть gauge-вертикальной.

\section{Решающий цветовой тест}

Среди шести обязательных ненулевых блоков находятся
\begin{equation}
 \phi_{Q_LY_R}\in(3,1)_{2/3},qquad
 \phi_{X_Lu_R}\in(3,1)_{5/3}.
 \label{eq:v7-full-edge-selected-colored-fields}
\end{equation}
Фундаментальное представление $\mathbf3$ не содержит ненулевого
$SU(3)$-инвариантного вектора:
\begin{equation}
 (\mathbb C^3)^{SU(3)}=\{0\}.
 \label{eq:v7-full-edge-no-color-invariant-vector}
\end{equation}
Но вакуум \eqref{eq:v7-full-edge-weighted-vacuum} требует норму $\mu>0$
для обоих цветных полей. Поэтому его стабилизатор является собственной
подгруппой $SU(3)$; полный цвет неизбежно нарушен. Никакой положительный
вес $w_e$, семейная кратность или Real-полуслед этого вывода не меняют.

Калибровочные индексы рёбер дополнительно подтверждают неоднородность:
\begin{equation}
 (I_1,I_2,I_3)_{E_*}
 =\left(\frac{29}{3},0,1\right),qquad
 (I_1,I_2,I_3)_{E\setminus E_*}
 =\left(\frac{10}{3},2,\frac12\right).
 \label{eq:v7-full-edge-sector-indices}
\end{equation}
На всех одиннадцати блоках получается
\begin{equation}
 (I_1,I_2,I_3)_{E}=\left(13,2,\frac32\right),
 \label{eq:v7-full-edge-total-indices}
\end{equation}
так что один прежний индекс $q_G=2$ не описывает полный carrier.

\begin{theorem}[Устойчивость опоры и цветовой запрет]
Полный минимальный gauge-подъём сохраняет точную опору $P_*$ при любых
положительных следовых весах: шесть выбранных блоков остаются ненулевыми,
пять конкурирующих --- нулевыми. Однако два выбранных блока являются
цветовыми триплетами, а Hodge-вакуум требует их ненулевых норм. Поэтому
текущий шестирёберный вакуум не сохраняет $SU(3)_c$ и не может быть
физическим вакуумом Стандартной модели.
\end{theorem}

\begin{nogocriterion}[Запрет цветного фундаментального конденсата]
Нельзя скрыть цветовой разрыв нормировкой trace, gauge-quotient или выбором
ориентации триплетов. Полный $SU(3)_c$ сохраняется только при нулевых
фундаментальных цветных вакуумных значениях. Рёбра $Q_LY_R$ и $X_Lu_R$
могут остаться переходными возбуждениями или входить в составной
цветосинглетный оператор, но не могут одновременно удовлетворять
$\|\phi_e\|=\mu>0$ в физическом вакууме.
\end{nogocriterion}

\begin{protocol}
\begin{enumerate}
 \item полных минимальных multiplet-блоков: \textbf{$11$};
 \item выбранная комплексная размерность: \textbf{$10$};
 \item конкурирующая комплексная размерность: \textbf{$11$};
 \item опора $6/11$ при положительных весах: \textbf{сохраняется};
 \item сигнатура нуля: \textbf{$(20,0,22)$};
 \item сигнатура вакуума до quotient: \textbf{$(0,14,28)$};
 \item выбранных цветных триплетов: \textbf{$2$};
 \item ненулевой $SU(3)$-инвариантный вектор в $\mathbf3$: \textbf{нет};
 \item сохранение цвета вакуумом: \textbf{нет};
 \item общий индекс рёбер: \textbf{$(13,2,3/2)$};
 \item численное gauge-замыкание: \textbf{не достигнуто};
 \item итог: \textbf{структурный селектор устойчив, физический вакуум закрыт};
 \item следующий гейт: \textbf{цветосохраняющий составной циклический
 родитель}.
\end{enumerate}
\end{protocol}

Машинный сертификат сохранён в
\path{s2t_v7_full_gauge_weighted_edge_carrier_gate_results.json}.

\begin{thebibliography}{9}
\bibitem{v7-full-edge-spectral-action}
A.~H.~Chamseddine, A.~Connes,
\emph{The Spectral Action Principle},
Communications in Mathematical Physics 186 (1997), 731--750;
arXiv:hep-th/9606001.
\bibitem{v7-full-edge-krajewski}
T.~Krajewski,
\emph{Classification of Finite Spectral Triples},
Journal of Geometry and Physics 28 (1998), 1--30;
arXiv:hep-th/9701081.
\end{thebibliography}